\section{ПРЕДМЕТНАЯ ОБЛАСТЬ И ПОСТАНОВКА ЗАДАЧИ}

\subsection{Рисуночные тесты как инструмент психологической диагностики}

Рисуночные тесты относятся к классу проективных психодиагностических методик. Их основная идея заключается в том, что в свободном изобразительном продукте отражаются устойчивые личностные особенности, эмоциональные состояния и социальные установки автора рисунка~\cite{venger}. Среди множества методик особое место занимает тест <<Рисунок человека>>: ребёнку предлагается нарисовать человека на чистом листе бумаги, после чего психолог анализирует получившееся изображение по большому числу формальных и содержательных критериев.

Методология интерпретации, разработанная А.~Л.~Венгером~\cite{venger}, применяется на практике российскими специалистами. Каждый рисунок описывается по более чем сотне критериев, объединённых в смысловые блоки: расположение на листе, размер рисунка, нажим, тело, ноги, одежда и другие. Каждой комбинации значений сопоставлен набор гипотетических психологических интерпретаций. Существенно, что одна и та же черта обычно поддерживается несколькими признаками рисунка одновременно, поэтому при интерпретации специалист строит общую картину, учитывая количество подтверждающих признаков.

\subsection{Список психологических черт и канонические признаки рисунка}

В настоящей работе использован экспертно составленный список из 26 канонических психологических черт --- тревожность личностная, состояние тревоги, демонстративность, агрессивность и другие --- каждая из которых снабжена кратким текстовым описанием. Список из 117 наблюдаемых признаков рисунка предоставлен научным руководителем. Признаки разделены на бинарные и категориальные.

\subsection{Современные подходы к классификации изображений}

Задача распознавания признаков на изображении сводится к задаче классификации: для каждого бинарного признака требуется двухклассовый классификатор, для каждого категориального --- мультиклассовый. С учётом совместного множества признаков задача относится к классу multi-label / multi-task classification.

Современным стандартом решения подобных задач являются свёрточные нейронные сети~\cite{imagenet}. Одним из наиболее распространённых семейств таких архитектур является ResNet~\cite{resnet}, в котором используются остаточные соединения, обеспечивающие устойчивое обучение глубоких сетей. Для multi-task обучения распространён подход с общим бэкбоном и отдельными классификационными головами для каждой задачи: совместное обучение позволяет частично разделять пространство признаков между задачами и сокращать требуемый объём данных, однако уязвимо к конкуренции градиентов. Для смягчения этого эффекта применяются взвешивание функции потерь и стратификация выборки.

Помимо собственно задачи распознавания, дополнительным аспектом разработки является способ предоставления результата конечному пользователю. Современным стандартом доставки функциональности машинного обучения в виде сервиса является микросервисная архитектура: каждый функциональный модуль системы оформляется как независимый процесс, общающийся с соседними по протоколу HTTP, а вся совокупность модулей развёртывается через систему контейнеризации. Такой подход обеспечивает технологическую независимость модулей, отказоустойчивость, простоту масштабирования и развёртывания.

\subsection{Особенности задачи и подход к её решению}

Особенностями детских рисунков являются малый контраст линий относительно фона, разнообразие стилей и плотности изображения, размещение фигуры в произвольной области листа, наличие посторонних деталей. Эти особенности требуют этапа предобработки --- локализации фигуры на странице и извлечения её ограничивающего прямоугольника. В работе для этой задачи разработан собственный алгоритм на основе классических операций обработки изображений, что обеспечивает прозрачность и воспроизводимость поведения модуля и не требует внешних предобученных моделей детекции.

После кропа фигуры и приведения её к фиксированному квадратному размеру теряется информация о масштабе и положении фигуры на исходном листе, существенная для ряда признаков, например: смещение. Для сохранения этой информации в работе используется двухпотоковая архитектура: помимо самого изображения, в модель подаётся дополнительный вектор геометрических метаданных bounding box. Подробное описание архитектуры приводится в Главе~2.

\subsection{Постановка задачи}

В рамках данной работы решается следующая задача. Дан датасет из 8179 изображений детских рисунков формата png/jpeg, сопровождаемый таблицей разметки в формате CSV. Каждой записи сопоставлены значения 117 признаков, проставленные экспертами-психологами по методологии А.~Л.~Венгера. База изображений и разметка предоставлены Московским государственным психолого-педагогическим университетом; рисунки получены из образовательных учреждений Самарской, Московской, Вологодской областей и Республики Чувашия.

Требуется:

\begin{itemize}
    \item разработать алгоритм предобработки изображений, обеспечивающий извлечение фигуры с произвольного листа;
    \item спроектировать и обучить нейросетевую модель, классифицирующую все 117 признаков по входному рисунку;
    \item обеспечить корректную обработку дисбаланса классов и сильной разнородности количеств положительных примеров между признаками;
    \item составить словарь сопоставления признаков с 26 каноническими психологическими чертами и реализовать на его основе модуль генерации текстового психологического портрета;
    \item провести экспериментальное исследование качества модели на тестовой выборке;
    \item для удобства практического использования оформить полученные модули в виде микросервисного веб-сервиса, доступного через браузер.
\end{itemize}

Целевые показатели качества по тестовой выборке --- F1-мера (micro), F1-мера (weighted) и F1-мера (macro). Все три варианта усреднения приводятся параллельно: F1-мера (micro) и F1-мера (weighted) отражают долю верных предсказаний на типичных рисунках, F1-мера (macro) --- качество распознавания всех классов одновременно, в том числе минорных.
